\chapter{Transformata wavelet și algoritmul lui Mallat}
\label{ch:wavelets}

\section{De ce wavelets}

La un cutremur, unda P sosește la $t=12$s și unda S la $t=38$s. Un spectru
Fourier calculat pe tot semnalul nu poate spune când s-a schimbat regimul,
ceea ce lipsește e localizarea. Wavelets rezolvă exact asta: în loc de o
sinusoidă infinită, folosesc un „wavelet mamă” $\psi$, o undă scurtă, cu
energie concentrată local, scalată și translatată,
\[
\psi_{a,b}(t) = \frac{1}{\sqrt{|a|}}\,\psi\!\left(\frac{t-b}{a}\right),
\]
unde $b$ e momentul de timp și $a$ scala, invers proporțională cu
frecvența, deci $a$ mic prinde detalii de frecvență înaltă, $a$ mare prinde
structuri de frecvență joasă. Rezultatul e o analiză multi-rezoluție:
semnale nestaționare (muzică, ECG, cutremure) capătă o descriere în care
timpul și frecvența apar simultan.

Compromisul e formalizat de principiul de incertitudine al lui Heisenberg,
$\Delta t \cdot \Delta f \ge 1/(4\pi)$~\cite{mallat2008}: fiecare metodă
împarte planul timp-frecvență în dreptunghiuri de arie constantă, dar forma
lor diferă. Fourier pur dă rezoluție infinită în frecvență și nulă în timp.
STFT, cu fereastră fixă, dă un compromis identic la orice frecvență.
Wavelets adaptează fereastra: îngustă în timp la frecvențe înalte, îngustă
în frecvență la frecvențe joase (Figura~\ref{fig:heisenberg-scalogram},
stânga), forma care se potrivește cel mai bine cu semnale reale, unde
tranzițiile scurte (un atac de tobă, un complex QRS) coexistă cu structuri
lungi (un bas continuu, linia de bază a unui ECG).

Scalograma (Figura~\ref{fig:heisenberg-scalogram}, dreapta) arată aceeași
idee pe un semnal concret: transformata wavelet continuă a unui chirp
modulat în amplitudine produce o bandă diagonală, frecvența crește cu
timpul, a cărei intensitate urmărește chiar amplitudinea semnalului de
intrare.

\begin{figure}[htbp]
\centering
\begin{subfigure}[t]{0.485\textwidth}
\includegraphics[width=\textwidth]{22-heisenberg-boxes.png}
\caption{Cutii de arie constantă: înguste în timp la frecvențe înalte, înguste
  în frecvență la frecvențe joase.}
\end{subfigure}
\hfill
\begin{subfigure}[t]{0.485\textwidth}
\includegraphics[width=\textwidth]{23-scalogram.png}
\caption{Scalograma unui chirp modulat: banda diagonală urcă în frecvență și
  respiră odată cu amplitudinea.}
\end{subfigure}
\caption{Compromisul timp-frecvență, teoretic (stânga) și pe un semnal
  concret (dreapta).}
\label{fig:heisenberg-scalogram}
\end{figure}

\section{Familii de wavelets}

Nu orice undă scurtă e un wavelet valid. Condiția de admisibilitate cere ca
$\int_0^\infty |\hat\psi(\omega)|^2/\omega\,d\omega$ să fie finită, ceea ce
pentru un $\psi$ integrabil forțează media lui nulă,
$\int \psi(t)\,dt = 0$~\cite{daubechies1992}, de-aici numele de familie: un
wavelet oscilează în jurul lui zero, nu se așază pe o valoare. Funcția de
scalare $\phi$ și wavelet-ul $\psi$ se leagă recursiv prin coeficienții
filtrului,
\[
\phi(t) = \sqrt2\sum_k h_k\,\phi(2t-k), \qquad
\psi(t) = \sqrt2\sum_k g_k\,\phi(2t-k),
\]
cu $g$ obținut din $h$ prin aceeași relație de filtre-oglindă-conjugate din
Secțiunea~\ref{sec:filtre}~\cite{estebangaland1977,smithbarnwell1984}.
Familia Daubechies~\cite{daubechies1988}, Symlet, Coiflet, Morlet și altele
diferă prin numărul de coeficienți și prin cât de simetrice sunt, toate
schimbă doar $\phi$ și $\psi$, algoritmul de mai jos rămâne același.

Transformata continuă, $W(a,b) = \int f(t)\cdot\frac{1}{\sqrt a}\psi^*
\left(\frac{t-b}{a}\right)dt$, evaluează $(a,b)$ pe un continuu, deci
informația se repetă. Transformata discretă eșantionează diadic, $a=2^j$,
$b=k\cdot 2^j$, și produce exact atâția coeficienți câți sunt necesari
pentru o bază ortonormală, fără redundanță~\cite{mallat2008}, baza pe care
lucrează algoritmul lui Mallat, mai jos. Un demo din tur desenează live
forma mamă pentru orice familie disponibilă în PyWavelets, citind direct
rezultatul funcției \txt{wavefun()} a bibliotecii. O sinusoidă simplă e
inclusă explicit ca ne-exemplu, fiindcă n-are nicio localizare în timp, deci
nu e wavelet.

\begin{figure}[htbp]
\centering
\includegraphics[width=\textwidth]{19-wavelet-families.png}
\caption{Cele patru rezultate din spatele familiilor wavelet: media nulă,
  ecuațiile de scalare/wavelet, perechea CQF și non-redundanța DWT, fiecare cu
  citarea lui.}
\label{fig:wavelet-families}
\end{figure}

\section{Wavelet-ul complex și faza}
\label{sec:complex}

Familiile de până acum sunt reale: un coeficient mare poate să însemne
rezonanță puternică, dar semnul lui depinde de unde anume a căzut wavelet-ul
peste oscilație. Mută semnalul cu o fracțiune de perioadă și coeficientul se
schimbă, deși conținutul nu s-a schimbat deloc. Wavelet-ul Morlet complex
rezolvă exact asta, pentru că e o gaussiană înmulțită cu o exponențială
complexă,
\[
  \psi(t) = e^{-t^2/2} \cdot e^{i\omega t}
          = e^{-t^2/2}\left(\cos(\omega t) + i \sin(\omega t)\right),
\]
adică două wavelet-uri reale defazate cu un sfert de perioadă, folosite
împreună. Coeficientul devine un număr complex, din care se citesc separat
magnitudinea, care spune \emph{cât} de puternică e rezonanța și nu se mai
schimbă la deplasarea semnalului, și faza, care spune \emph{unde} a căzut
oscilația în interiorul ferestrei.

Ecranul din tur (Figura~\ref{fig:complex-wavelet}) arată wavelet-ul ca o
spirală în trei dimensiuni: timpul pe o axă, partea reală pe a doua, partea
imaginară pe a treia, cu culoarea codificând faza de la $-\pi$ la $\pi$.
Proiecția pe fiecare perete dă cele două componente reale, iar spirala e
motivul pentru care scalograma din secțiunea anterioară arată benzi netede și
nu o alternanță de dungi luminoase și întunecate: acolo se desenează
magnitudinea coeficientului complex, care nu depinde de fază.

Aceeași proprietate e motivul pentru care wavelet-urile complexe se folosesc
la detectarea trăsăturilor invariante la translație și la analiza mișcării,
unde contează ce s-a întâmplat, nu la ce microsecundă a căzut fereastra.

\begin{figure}[htbp]
\centering
\includegraphics[width=\textwidth]{24-complex-wavelet.png}
\caption{Wavelet-ul Morlet complex ca spirală: partea reală proiectată pe un
  perete, partea imaginară pe celălalt, iar culoarea curbei centrale codifică
  faza. Frecvența și unghiul de rotație sunt reglabile din dreapta.}
\label{fig:complex-wavelet}
\end{figure}

\section{Algoritmul lui Mallat}
\label{sec:mallat}

Formal, un semnal se descompune prin proiecție pe funcția de scalare și pe
wavelet la nivelul $j$,
\[
c_{j,k} = \int x(t)\,\phi_{j,k}(t)\,dt, \qquad d_{j,k} = \int x(t)\,\psi_{j,k}(t)\,dt.
\]
Contribuția lui Mallat e observația care face aceste integrale inutile de
calculat direct: coeficienții unui nivel se obțin din nivelul anterior doar
prin filtrare și decimare cu doi,
\[
c_{j+1,k} = \sum_n h_0[n-2k]\,c_{j,n}, \qquad
d_{j+1,k} = \sum_n h_1[n-2k]\,c_{j,n},
\]
unde $h_0$ e filtrul de aproximare $h$ și $h_1$ filtrul de detaliu $g$ din
ecuațiile de mai sus, scrise acum cu indici fiindcă apar împreună,
fără nicio integrală calculată efectiv~\cite{mallat1989,mallat2008}. E o
recursie exactă: fiecare nivel iese din cel anterior printr-o singură
trecere de filtrare, deci descompunerea pe $L$ niveluri costă $L$ perechi de
filtrări, nu $L$ integrale separate. În cod, aceeași recursie e
\txt{pywt.wavedec2(imagine, wavelet, level=niveluri)}, care întoarce
aproximarea cea mai grosieră plus câte un triplet de subbenzi de detaliu pe
nivel. Reconstrucția, \txt{pywt.waverec2}, trece coeficienții prin filtrele
de sinteză și îi însumează, operația inversă a celei de mai sus.

Pentru wavelet-ul Haar, singurul cu filtre de doi coeficienți, suma se scrie
pe un singur eșantion și devine o pereche de medie și diferență,
\[
x[2k] = \frac{cA[k]+cD[k]}{\sqrt2}, \qquad x[2k+1] = \frac{cA[k]-cD[k]}{\sqrt2},
\]
forma folosită de demo-ul de reconstrucție din tur. Pentru orice filtru mai
lung, un Daubechies sau biortogonalul 9/7 din capitolul de compresie, un
eșantion reconstruit primește contribuții de la mai multe perechi
$(cA,cD)$ vecine, iar formula de mai sus nu se mai aplică punct cu punct.

La imagini, aceeași pereche de filtre se aplică o dată pe rânduri și o dată
pe coloane, dând patru subbenzi, $LL$ (aproximare), $LH$/$HL$ (detalii
orizontale/verticale) și $HH$ (detalii diagonale). Analogia cu hărțile ajută:
la zoom mic vezi doar $LL$, structura globală. La zoom mare apar
$LH$/$HL$/$HH$, detaliile fine. Pasul se aplică recursiv doar pe $LL$.
High-pass pe coloane dă detalii \emph{orizontale} pentru că răspunde la
variația pe verticală, exact ce înseamnă o muchie orizontală.

Demo-ul educațional 1D (Figura~\ref{fig:mallat-1d-edu}) face vizibil un
nivel al algoritmului, pas cu pas, pe un semnal de 16 eșantioane, treaptă,
rampă, sinusoidă, puls sau o „melodie” de note discrete. Fereastra
filtrului, produsul și suma sunt animate explicit, apoi decimarea cu doi.
Interpretarea diferă cu forma semnalului, o treaptă produce un vârf de
detaliu exact la muchie, o melodie păstrează notele în aproximare și
marchează atacul fiecărei note în detaliu. Un glosar de notație, repetat pe
toate ecranele acestei secțiuni, ține evidența faptului că aceeași împărțire
circulă sub nume diferite în literatură și în cod: $c = cA = L = A$ pentru
aproximare, $d = cD = H = D$ pentru detaliu.

\begin{figure}[htbp]
\centering
\includegraphics[width=\textwidth]{28-mallat-1d-edu.png}
\caption{Un nivel al algoritmului lui Mallat pas cu pas: fereastra filtrului,
  produsul, suma și decimarea cu doi, pe un semnal treaptă de 16 eșantioane.}
\label{fig:mallat-1d-edu}
\end{figure}

Aplicat recursiv doar pe aproximare, pe patru niveluri
(Figura~\ref{fig:pyramid-recon}, stânga), pasul arată unde se duce energia
unui semnal real: peste 94\% rămâne concentrată în aproximarea cea mai
grosieră, A4. E numărul concret din spatele afirmației că „majoritatea
energiei unei imagini încape într-un număr mic de coeficienți”, premisa pe
care se sprijină toată povestea compresiei din Capitolul~\ref{ch:compresie}.
Un ultim ecran (Figura~\ref{fig:pyramid-recon}, dreapta) arată ce înseamnă
„reconstrucție perfectă” în aritmetică reală: dacă filtrele satisfac
condiția CQF, eroarea rămâne la precizia mașinii de calcul, iar PSNR-ul
raportat reflectă exact acea precizie finită, de aici pornește atât
compresia \emph{lossless}, cât și orice pierdere ulterioară, controlată prin
cuantizare.

\begin{figure}[htbp]
\centering
\begin{subfigure}[t]{0.485\textwidth}
\includegraphics[width=\textwidth]{31-pyramid-decomp.png}
\caption{Patru niveluri recursive: peste 94\% din energie rămâne în A4.}
\end{subfigure}
\hfill
\begin{subfigure}[t]{0.485\textwidth}
\includegraphics[width=\textwidth]{32-reconstruction.png}
\caption{Reconstrucție perfectă: eroarea rămâne la precizia mașinii, nu la
  zero exact.}
\end{subfigure}
\caption{Multi-nivel (stânga) și reconstrucția care închide bucla (dreapta).}
\label{fig:pyramid-recon}
\end{figure}

Pentru compunerea vizuală a subbenzilor într-un singur mozaic, cum apare pe
mai multe ecrane din tur, fiecare pereche de subbenzi e decupată la
dimensiunea minimă comună înainte de alăturare. PyWavelets poate produce
subbenzi care diferă cu un pixel pe dimensiuni impare de intrare.
